% Generated draft: requires mathematical and source review.
\documentclass[11pt]{article}
\usepackage[margin=25mm]{geometry}
\usepackage{fontspec}
\setmainfont{Latin Modern Roman}
\usepackage{amsmath,amssymb,mathtools,unicode-math}
\setmathfont{Latin Modern Math}
\usepackage{enumitem,needspace,xurl,hyperref}
\setlength{\parindent}{0pt}
\setlength{\parskip}{5pt}
\setlength{\emergencystretch}{3em}
\newcommand{\N}{\mathbb{N}}
\newcommand{\Z}{\mathbb{Z}}
\newcommand{\Q}{\mathbb{Q}}
\newcommand{\R}{\mathbb{R}}
\newcommand{\C}{\mathbb{C}}
\newcommand{\F}{\mathbb{F}}
\newcommand{\A}{\mathbb{A}}
\newcommand{\PP}{\mathbb{P}}
\newcommand{\E}{\mathbb{E}}
\newcommand{\eps}{\varepsilon}
\newcommand{\dd}{\,\mathrm{d}}
\begin{document}
\section*{UnsolvedMath ID: 30002480}
\textbf{Source problem code:} catalogue entry
\subsection{Definitions and mathematical statement}
Regard cardinals as initial ordinals. For an infinite cardinal $\lambda$, its cofinality $\operatorname{cf}(\lambda)$ is the least cardinality of an unbounded subset of $\lambda$. The cardinal $\lambda$ is singular if $\operatorname{cf}(\lambda)<\lambda$. Write $[\kappa]^2$ for the set of two-element subsets of $\kappa$.

The supplied record does not define $\operatorname{Pr}_1$. One explicit candidate interpretation of $\operatorname{Pr}_1(\kappa,\mu,\tau,\theta)$ is the existence of a coloring
\[
c:[\kappa]^2\longrightarrow\tau
\]
such that, for every family $\mathcal A$ of $\mu$ pairwise disjoint nonempty subsets of $\kappa$, each of cardinality less than $\theta$, and every $\xi<\tau$, there exist distinct $a,b\in\mathcal A$ satisfying
\[
\forall x\in a\ \forall y\in b,\qquad c(\{x,y\})=\xi.
\]
This is an explicit interpretation for review, not a definition verified against the supplied references. In particular, the record does not resolve whether the intended definition permits varying sizes below $\theta$ or instead quantifies separately over families of sets of one fixed size $\sigma<\theta$.

With the intended source convention for $\operatorname{Pr}_1$ still to be confirmed, the recorded questions are whether
\[
\forall\lambda\;\bigl[\lambda\text{ is an infinite singular cardinal}\ \Longrightarrow\ \operatorname{Pr}_1(\lambda,\lambda,\lambda,\operatorname{cf}(\lambda))\bigr]
\]
holds, and, as a weaker question, whether
\[
\forall\lambda\;\bigl[\lambda\text{ is an infinite singular cardinal}\ \Longrightarrow\ \operatorname{Pr}_1(\lambda,\lambda,\lambda,2)\bigr]
\]
holds. Under the candidate interpretation above, the parameter $2$ restricts the nonempty members of $\mathcal A$ to singletons.

\textit{Unresolved source qualification.} All displayed arguments are $\lambda$ exactly as supplied. One supplied literature reference is labelled as concerning successors of singular cardinals; the record does not explain its relationship to these questions at singular cardinals themselves. No replacement of $\lambda$ by $\lambda^+$ is made.
\subsection{Short English statement}
The source asks whether every singular cardinal admits the stated strong coloring, first for disjoint sets of size below its cofinality and then for singletons. This explanation uses the explicit candidate definition above; the intended convention for $\operatorname{Pr}_1$ and the connection with the supplied reference about successors require review.
\subsection{Sources}
\begin{itemize}
\item Supplied local UnsolvedMath registry record, permanent ID $30002480$, title \emph{Strong Colorings at Singular Cardinals}, collection \emph{Oberwolfach Reports Open Problems}, exact source code \texttt{catalogue entry}. The recorded proposed year is 2014; no individual proposer is supplied.
\item Supplied source citation: \emph{Set Theory} (2014), Oberwolfach Reports, DOI 10.4171/owr/2014/02, pp.\ 91--144. \url{https://doi.org/10.4171/owr/2014/02}. The record describes these reports as documenting mathematical workshops and their open-problem discussions.
\item Supplied literature source, labelled ``Oberwolfach set theory report'': \url{https://ems.press/content/serial-article-files/46488}.
\item Supplied literature source, labelled ``Shelah, Pr1 at successors of singulars'': \url{https://intlpress.com/site/pub/files/_fulltext/journals/mrl/2014/0021/0006/MRL-2014-0021-0006-a009.pdf}.
\item Literature checked at: \texttt{2026-08-21}. This is supplied dated curation, not an independently verified assessment. It labels the problem open and reports that the stronger relation for every singular cardinal is described in current primary literature as the oldest open problem in this line of work. The background characterizes this as dated literature triage derived from the supplied OWR curation and requests independent verification of time-sensitive status and cited claims.
\item The supplied statement status is \texttt{corrected\_verified}; its accompanying note reports manual semantic, typographic, and LaTeX review of a meaning-preserving cleanup. This is a supplied review claim, not independent mathematical verification of this draft.
\end{itemize}
\end{document}
